\scanpage{124}
% line head | text | 73
\mspage{73}
% line 1 | text | On a sûrement un foncteur
On a sûrement un foncteur
% line 2 | display | Comparison functor
\[
{}_{u'}\backslash^{U}_{/t}\longrightarrow
\bigl({}_{u'}\backslash^{\unsure{U'|U}}\times(T_{/t})^{\mathrm o}\bigr)_{/H}
\]
% line 3 | text | qui associe à un triple
qui associe à un triple
% line 4 | display | (u, alpha, beta prime)
\[
(u,\alpha,\beta'):\quad u\in\operatorname{Ob}U,\quad\alpha:u\unsure{-|\to}t,\quad
\beta':u'\to f_U(u)
\]
% line 5–7 | text | le triple
le triple $(\bar u,\bar t,\unsure{\xi|\zeta})$ ($\xi\in H(\bar u,\bar t)$) \deleted{du \unsure{sec.}} dans
le second membre, avec $\bar u$ dans $\mathord{}_{u'}\backslash^{U'}$, $\bar t$ dans $T_{/t}$,
défini ainsi :
% line 8 | display | barred u, barred t, xi
\[
\bar u=(u,\beta'),\quad \bar t=(t,\mathrm{id}_t),\quad
\underset{\substack{\shortparallel\\\alpha}}{\xi}\in H(\bar u,\bar t)=\operatorname{Hom}(u,t).
\]

% line 9 | text | En somme
En somme,
% line 10 | display | Fibre at the initial object
\[
\begin{aligned}
{}_{u'}\backslash^{U}_{/t}&=\text{fibre de }\bigl({}_{u'}\backslash^{U}\times(T_{/t})^{\mathrm o}\bigr)_{/\unsure{H|t}}
\text{ en l'objet \deleted{final}\inserted{initial} }\bar t\\
&\hspace{17em}\text{de }(T_{/t})^{\unsure{\mathrm o|'}}
\end{aligned}
\]
% line 11 | text | par la projection naturelle
par la projection naturelle, flèche en
% line 12–14 | diagram | three levels of the natural projection
\diagram{\begin{tikzcd}[ampersand replacement=\&,row sep=small]
\bigl({}_{u'}\backslash^{U}\times(T_{/t})^{\mathrm o}\bigr)_{/H}\arrow[d]\\
{}_{u'}\backslash^{U}\times(T_{/t})^{\mathrm o}\arrow[d,"\pi"]\\
(T_{/t})^{\mathrm o}
\end{tikzcd}}
% line 15–21 | text | Il faut voir que l'inclusion
Il faut voir que l'inclusion de la fibre donne
une bijection sur les $\pi_0$, ou ce qui revient
au $=$, que l'inclusion de la fibre
$\mathord{}_{u'}\backslash^{U}$ \uline{de $\pi$} en l'objet initial $(t,\mathrm{id}_t)$ de $(T_{/t})^{\mathrm o}$,
donne une bijection sur les $\lim_{\rightarrow}$ \unsure{du|des}
\unsure{foncteur|foncteurs} $H$ !\unsure{!} Donc l'inclusion a intérêt
à être asphérique, ou du moins co-cofinal.
% archive sheet number: 124
